\begin{textsample}{POS Dim 5 – global\_north\_2023\_11 – Score 18.00 – global\_north\_2023\_11/103468397.txt}  \label{ex:f5_pos_411}
Urgent Need For New Zealand ' s Humanitarian Voice Rt Hon Chris Hipkins Prime Minister ( Caretaker ) \&amp ; Christopher Luxon Prime Minister ( Elect ) Parliament Wellington Further to our previous letter regarding this matter, Aotearoa New Zealand has always been an \textbf{international} beacon for timely humanitarian and peaceful initiatives - whether it was our nuclear free position in the Pacific, or our stance against participation in the Iraq war without UN \textbf{sanction}, or our nation ' s response to the tragic terror attacks of 15 March 2019.
Aotearoa New Zealand has been \textbf{recognised} globally as being in the forefront of the \textbf{international} moral conscience at times when there was a need for a peaceful pathway forward.
We write now to once again raise our grave concerns and strongest plea for a Government response to the ever-worsening situation in Gaza.
We note with appreciation New Zealand ' s recent voting at the UN General Assembly on Gaza.
In keeping with the above we request the Government ' s urgent diplomatic advocacy at both the UN and other global forums in seeking a complete cessation @ @ @ @ @ @ @ @ @ @ the tragic truths and harsh reality of over 10,000 civilians killed by the Israeli military, including over 4,600 children.
Of particular and unconscionable note are the harrowing images of rows of neonatal infants lying on the cold concrete floor.
Such \textbf{acts} of deliberate and wanton \textbf{genocide} have not been seen in recent history and this is an appropriate time for the New Zealand Government to raise its voice.
Advertisement - \textbf{scroll} to continue reading Are you getting our free newsletter?
We urge the Government, on the basis of every moral and humanitarian \textbf{principle}, to \textbf{unequivocally} condemn such \textbf{genocidal} assaults and wholescale bombing of innocent civilians as being contrary to all \textbf{international} norms of civil behaviour and law.
There are four specific positions that we feel New Zealand should immediately adopt as peaceful but profound initiatives in response to the above : As a sign of New Zealand ' s concern for the escalating humanitarian crisis, and the breaches of \textbf{international} law, the Israeli Ambassador to New Zealand should be expelled.
Particularly noting the Ambassador ' s @ @ @ @ @ @ @ @ @ @ the past.
That New Zealand suspend its March 2020 agreement with the Government of Israel on \textbf{cooperation} in technological innovation research and development as it may have the potential for military application.
To immediately stop Ministry of Business, Innovation, and Employment ( MBIE ) using the Cobweb Technologies software which was developed by ex-Israeli intelligence officials and has the potential to track, trace and monitor the private messages of Palestinians in New Zealand with their families and friends in Gaza.
It should be noted that unlike other New Zealand intelligence-gathering agencies, there is no external \textbf{independent} monitoring of the usage of this software by MBIE.
This lack of democratic oversight is most alarming.
The immediate increase of humanitarian aid to \textbf{international} agencies assisting with relief efforts in Gaza, including WHO, Red Crescent and UNRWA.
This plea is based on New Zealand ' s principled stance as we have recently taken.
We note the parallels between the New Zealand response to the situation in the Ukraine and that in Gaza.
The New Zealand response @ @ @ @ @ @ @ @ @ @ sense of urgency, timeliness and appropriateness with over \$83 million \textbf{committed}.
The current response to Gaza has been neither timely nor commensurate to the gravity of the situation.
As we have been reminded by the United Nations,"The situation in the Gaza Strip is a growing stain on our collective conscience."History will judge us for our lack of decisive action if we do not take a principled stance at this time.
Nga mihi, Ibrar Sheikh President, FIANZ Did you know \textbf{Scoop} has an Ethical Paywall?
If you ' re using \textbf{Scoop} for work, your organisation needs to pay a small \textbf{license} \textbf{fee} with \textbf{Scoop} \textbf{Pro}.
We think that ' s \textbf{fair}, because your organisation is benefiting from using our news resources.
In return, we ' ll also give your team access to \textbf{pro} news tools and keep \textbf{Scoop} free for personal use, because public access to news is important!
Among the key issues facing diplomats is securing the release of a reported 199 Israeli hostages, @ @ @ @ @ @ @ @ @ @,"says Emergency Relief Coordinator Martin Griffiths."This war was started by taking those hostages.
Of course, there ' s a history between Palestinian people and the Israeli people, and I ' m not denying any of that.
But that \textbf{act} alone lit a fire, which can only be put out with the release of those hostages."More Families in western Afghanistan are reeling after a fourth earthquake hit Herat Province, crumbling buildings and forcing people to flee once again, with thousands now living in tents exposed to fierce winds and dust storms.
The latest 6.3 magnitude earthquake hit 30 km outside of Herat on Sunday, shattering communities still reeling from strong and shallow aftershocks.
More UN Spokesperson St? phane Dujarric said some 1.1M people would be expected to leave northern Gaza and that such a movement would be"impossible"without devastating humanitarian consequences and appeals for the order to be rescinded.
The WHO joined the call for Israel to rescind the relocation order, which amounted to a"death @ @ @ @ @ @ @ @ @ @, reports indicated that fixed-line internet, mobile data, SMS, telephone, and TV networks are all seriously compromised.
With significant and increasing damage to the electrical grid, orders by the Israeli Ministry of Energy to stop supplying electricity and the last remaining power station now out of fuel, many are no longer able to charge devices that are \textbf{essential} to communicate and access information.
More

% matched lemmas: act, commit, cooperation, essential, fair, fee, genocidal, genocide, independent, international, license, principle, pro, recognise, sanction, scoop, scroll, unequivocally
\end{textsample}
